\section{Билет 19. Уравнение плоской волны. Фазовая скорость. Уравнение плоской волны в среде. Уравнение плоской волны, распространяющейся в произвольном направлении.}

\begin{definition}
Рассмотрим источник гармонических колебаний, находящийся в начале координат ($x=0$). Его колебания описываются законом:
\begin{equation}
    \xi(0, t) = A \cos(\omega t).
\end{equation}
Волна распространяется вдоль оси $x$ со скоростью $v$. Колебания дойдут до точки с координатой $x$ с временной задержкой $\tau = x/v$. Поскольку среда однородна и поглощением пренебрегаем, амплитуда и частота сохраняются. Смещение частиц в точке $x$ в момент $t$ равно смещению источника в момент $t - \tau$:
\begin{equation}
    \xi(x, t) = A \cos\left[\omega \left(t - \frac{x}{v}\right)\right]. \label{eq:plane_wave_19}
\end{equation}
Вводя волновое число $k = \omega/v = 2\pi/\lambda$, получаем компактную запись уравнения плоской бегущей волны:
\begin{equation}
    \xi(x, t) = A \cos(\omega t - kx). \label{eq:plane_wave_k}
\end{equation}
\end{definition}

\begin{theorem}[Фазовая скорость]
\textbf{Фазовой скоростью} называется скорость перемещения точки с постоянной фазой колебаний. Пусть фаза волны $\Phi(x, t) = \omega t - kx$ сохраняет постоянное значение:
\begin{equation}
    \omega t - kx = \text{const}.
\end{equation}
Дифференцируя это тождество по времени, получаем:
\begin{equation}
    \omega \, dt - k \, dx = 0 \quad \Rightarrow \quad v_{\text{ф}} = \frac{dx}{dt} = \frac{\omega}{k}.
\end{equation}
Подставляя $k = \omega/v$, убеждаемся, что фазовая скорость совпадает со скоростью распространения волны в среде:
\begin{equation}
    v_{\text{ф}} = v.
\end{equation}
Для волны, бегущей в направлении убывания $x$, уравнение имеет вид $\xi = A\cos(\omega t + kx)$, а фазовая скорость $v_{\text{ф}} = -v$.
\end{theorem}

\begin{definition}[Уравнение плоской волны в поглощающей среде]
В реальной среде часть энергии волны расходуется на преодоление внутреннего трения и нагрев, поэтому амплитуда убывает с расстоянием. В однородной среде затухание происходит по экспоненциальному закону:
\begin{equation}
    A(x) = A_0 e^{-\gamma x},
\end{equation}
где $\gamma$ --- линейный коэффициент поглощения (зависит от свойств среды и частоты волны). Уравнение плоской волны в среде с поглощением принимает вид:
\begin{equation}
    \xi(x, t) = A_0 e^{-\gamma x} \cos(\omega t - kx). \label{eq:dissipative_wave}
\end{equation}
\end{definition}

\begin{definition}[Произвольное направление распространения]
Если волна распространяется в трёхмерном пространстве в произвольном направлении, заданном единичным вектором $\vec{n}$, то аргумент косинуса заменяется на скалярное произведение волнового вектора $\vec{k}$ и радиус-вектора $\vec{r}$ точки наблюдения:
\begin{equation}
    \xi(\vec{r}, t) = A \cos(\omega t - \vec{k} \cdot \vec{r}), \label{eq:3d_wave}
\end{equation}
где $\vec{k} = k \vec{n} = \frac{2\pi}{\lambda} \vec{n}$ --- \textbf{волновой вектор}. Его направление совпадает с направлением распространения волны, а модуль равен волновому числу $k$. Уравнение \eqref{eq:3d_wave} описывает плоскую монохроматическую волну, волновые поверхности которой являются плоскостями, перпендикулярными $\vec{k}$.
\end{definition}
